% Preface: the README's introduction, About the Course, During the Course, Learning Outcomes,
% Written Examinations and Building sections, and the prerequisites, course plan and software from
% info/README.md, as the front matter of a book.

\chapter{Preface}
\markboth{Preface}{Preface}

This book is for software and embedded engineers who know classical digital design and want to
write VHDL. It takes you from a gate network on paper to a design running on a real FPGA (a
\emph{field-programmable gate array}: a chip full of logic you configure into the circuit you want,
rather than one whose function was fixed when it was manufactured).

\textbf{The subject is VHDL.} Gates, truth tables, Boolean algebra and Karnaugh maps are refreshed
rather than taught: they appear where the language needs them, and in Chapter~\ref{c8:ch} where the
state machine is designed by hand before it is written.

The book covers writing synthesizable VHDL for combinational and sequential hardware:
metastability and synchronization, counters, timers and finite state machines, each reasoned about
as a circuit first and then expressed in the language, running on a real FPGA. Topics include:
\begin{itemize}
  \item Combinational logic: gates, truth tables, Boolean algebra, and VHDL from the first chapter.
  \item Larger networks: Karnaugh maps, multiplexers, \code{process} and \code{case}.
  \item Sequential logic: registers, D flip-flops, clocking, and edge detection.
  \item Metastability and synchronization: asynchronous inputs, synchronizers, debouncing.
  \item The VHDL language: signals, variables, subcomponents, and generics.
  \item Counters, shift registers, and timers.
  \item Finite state machines: designed by hand, then written in VHDL, Moore and Mealy.
\end{itemize}

\textbf{Wherever a circuit can usefully be drawn, it is drawn first}: reasoned about as gates, then
built and simulated by hand in CircuitVerse (\url{https://circuitverse.org/simulator}) before it is
ever written in VHDL. Every chapter that introduces a new kind of circuit draws one that way first.
The exception is Chapter~\ref{c5:ch}, which is about the language and has no new circuit in it; and
once a kind of circuit has been drawn, later designs of the same kind go straight into VHDL.

Verification works in three layers:
\begin{itemize}
  \item \textbf{By hand, in CircuitVerse}, before any VHDL exists. You derive what the circuit
    should do, build it, and then simulate it yourself against that prediction, exhaustively,
    since the circuits are small enough for that to be cheap. Nothing automated checks this step,
    deliberately: predicting a behaviour and then confirming it is the skill the rest of the book is
    built on.
  \item \textbf{In the exercises}, you write the VHDL and check it with a provided self-checking
    testbench (\file{<module>_tb.vhd}, runnable with GHDL). Nearly every VHDL exercise ships one,
    so nearly every exercise can be completed and verified on a laptop; the handful that do not say
    so, and say how to check by hand instead.
  \item \textbf{On the board}: the designs are synthesized in Quartus Prime Lite and run on a real
    Terasic DE0-CV board. In the course the book comes from, that is demonstrated from the front,
    and nobody needs a board of their own; Appendix~\ref{app:quartus} is the workflow, for a reader
    who has one.
\end{itemize}

\textbf{Writing} testbenches is not a topic here: this book only shows you how to run the ones
provided.

\textbf{Along the way you will build circuits, not just read about them.} You will:
\begin{itemize}
  \item Realize a logic function by hand as a gate network, and simulate it in CircuitVerse.
  \item Translate that same gate network into synthesizable VHDL.
  \item Build sequential circuits from D flip-flops: registers, counters, shift registers, and
    timers.
  \item Guard every asynchronous input with a double-flop synchronizer before it reaches synchronous
    logic.
  \item Design a Moore state machine by hand (diagram, table, Karnaugh maps, gate network), then
    write the same kind of machine in VHDL, and learn to recognize a Mealy machine and what it
    costs.
  \item Follow each design through Quartus Prime Lite and onto a Terasic DE0-CV FPGA board.
  \item Verify your own VHDL after each chapter by running the provided self-checking testbenches.
\end{itemize}
The examples and exercises are small, complete circuits, such as:
\begin{itemize}
  \item A gate network realizing a Boolean function, in CircuitVerse and in VHDL.
  \item A synchronized, edge-detected push button driving an LED.
  \item A free-running timer and a shift-register-driven LED pattern.
  \item A three-state Moore machine cycling an LED off, blinking, on, composed from a synchronizer
    and a timer.
  \item A two-state Mealy sequence detector, contrasted against the Moore machine it would take
    three states to write.
\end{itemize}
The book then closes on two exercises that are not small: a serial receiver and a serial
transmitter, each composed from modules you wrote in earlier chapters around a state machine you
design yourself, and each verified by a provided testbench.

\section*{What you should be able to do at the end}
Having worked through the book, you should be able to:
\begin{itemize}
  \item Realize a Boolean function as a gate network, and simulate it by hand in CircuitVerse.
  \item Read and write combinational and sequential VHDL: entities, architectures, signals,
    variables, and processes.
  \item Explain metastability and apply the double-flop synchronizer to any asynchronous input.
  \item Design and implement counters, shift registers, and timers in VHDL.
  \item Design a Moore state machine from a specification and implement it in VHDL, and recognize a
    Mealy machine on sight, knowing what its earlier output costs.
  \item Compose a design out of previously built modules (a reset synchronizer, a button
    synchronizer, a timer) rather than rewriting their logic.
  \item Explain how a VHDL design reaches an FPGA: synthesis, pin assignment, and programming the
    board in Quartus Prime Lite.
  \item Verify a VHDL design by running a self-checking testbench under GHDL.
\end{itemize}

\section*{What you are assumed to know}
The subject of this book is \textbf{VHDL}, not digital design. You are expected to arrive already
comfortable with:
\begin{itemize}
  \item Logic gates and truth tables, and reading a sum-of-products equation off a table.
  \item Boolean algebra: the basic identities and De Morgan's laws.
  \item Karnaugh maps, for functions of three or four variables.
  \item Registers, D flip-flops, and what a clock edge does.
  \item State machines as a concept: states, transitions, and the difference between Moore and
    Mealy.
\end{itemize}
None of that is taught from scratch, and this book is the poorer place to meet it first. The
junior-track course \emph{Digital Electronics}
(\url{https://github.com/qrtech-academy/digital-electronics}) teaches all five points above from the
ground up, derived by hand and simulated in CircuitVerse; this book is written to start where that
course ends, so read it first if any of the list is unfamiliar.

What is here instead is refreshed where the language needs it, briefly:
Chapter~\ref{c1:ch} restates gate and Boolean notation so the book has one convention,
Chapter~\ref{c2:ch} works a Karnaugh map end to end because minimization is what motivates the VHDL
that follows, and Chapter~\ref{c8:ch} designs one state machine by hand before writing it, because
seeing the state table become a \code{case} statement is the point of the chapter.

What the book does teach from nothing is the language and the toolchain: \code{entity} and
\code{architecture}, signals and variables, processes, generics, subcomponents, synthesis in
Quartus, and simulation with GHDL.

No programming language beyond general software fluency is assumed. Familiarity with C or a similar
language helps, since several explanations contrast VHDL's semantics with it.

\section*{How the book is organised}
The book is typeset from an eight-lecture course, and each chapter is one lecture. Its sections are
the parts of the lecture's appendix, in the same order and with the same subsections, so a chapter
here and a lecture directory in the repository can be read side by side:

\begin{booktable}{@{}p{5.2em}p{3.6em}L@{}}
  \thead{Chapter} & \thead{Lecture} & \thead{Topic} \\
  \midrule
  1 & \file{L01} & Combinational logic and first VHDL \\
  2 & \file{L02} & Larger networks, multiplexers and submodules \\
  3 & \file{L03} & Sequential logic \\
  4 & \file{L04} & Metastability and synchronization \\
  5 & \file{L05} & Variables and the hardware underneath \\
  6 & \file{L06} & Counters and shift registers \\
  7 & \file{L07} & Timers \\
  8 & \file{L08} & State machines \\
\end{booktable}

Every chapter opens with what is in it, works through the material, closes with a review of what
you should now be able to explain, and ends with its exercises. The appendices hold the short
answers to the exercises, the Quartus workflow, and the two written papers described below.

\section*{What the exercises look like}
The exercises are grouped into \textbf{sets} by topic, and each is labelled with the kinds of work
it asks for:

\begin{booktable}{@{}p{7.4em}L@{}}
  \thead{Kind} & \thead{What it asks of you} \\
  \midrule
  \kinds{Paper} & Work it by hand: a truth table, a derivation, a Karnaugh map, a trace through
    the clock edges, a calculation, or a question answered in a sentence or two. \\
  \kinds{CircuitVerse} & Build the circuit in CircuitVerse and simulate it, exhaustively, against
    the behaviour you predicted first. \\
  \kinds{VHDL} & Write a module, with the entity name and port order the exercise gives, and check
    it with the testbench it ships. \\
\end{booktable}

\textbf{The answers to the paper work are in Appendix~\ref{app:answers}}, short, and by exercise
number, so that an answer worked on paper can be checked on paper. \textbf{The VHDL is checked by
its testbench}, which is the point of shipping one, and reference solutions are published in the
companion repository, under \file{lectures/LNN/solutions}. Try each exercise before you read its
answer.

\section*{Two written papers, and what they are not}
Nothing in the course this book comes from is marked. Assessment is the exercises after every
chapter, the self-checking testbench nearly every one of them ships, and the circuit you predict and
then simulate by hand in CircuitVerse before any VHDL exists.

Appendices~\ref{exam:ch:about} to \ref{ansb:ch:answers} hold two three-hour papers with worked
solutions, and they check something else: \textbf{your own skills and knowledge, on paper, with
nothing in front of you.} Eight questions each, one per chapter, mixing theory with VHDL you write
out by hand and with designs that compile cleanly and are wrong anyway.

\textbf{They exist purely so that you can test yourself after working through the book.} They gate
nothing, they are not a qualification, and no part of the book requires them.

\textbf{Take one after the last chapter}, and after the serial receiver and transmitter that close
it. Both papers draw on all eight chapters, so sitting one partway through examines material you
have not read yet.

\section*{Setting up}
Everything in this book runs on an ordinary laptop, with no hardware of any kind. You need:
\begin{itemize}
  \item \textbf{CircuitVerse}, the free, browser-based logic simulator used to build and simulate
    every circuit by hand before it is written in VHDL, at
    \url{https://circuitverse.org/simulator}. Nothing to install.
  \item \textbf{GHDL}, the free, open-source VHDL analyzer and simulator, used to run the
    self-checking testbenches that verify the exercises. Running testbenches is covered in
    \secref{c2:app:c}.
  \item \textbf{Git} and \textbf{make}, to fetch the companion repository and run its checks.
\end{itemize}
On Linux, or on Windows through WSL with Ubuntu:

\begin{shell}
sudo apt -y update
sudo apt -y install git make ghdl
ghdl --version           # 3.x or newer; the repository's CI builds against 4.1.0
\end{shell}

\noindent Then clone the companion repository:

\begin{shell}
git clone https://github.com/qrtech-academy/digital-design-vhdl.git
\end{shell}

\noindent Every lecture directory holds its appendices, its worked examples, each with a reference
testbench, and an \file{exercises} directory with the testbench for every VHDL exercise, in a
directory named after the module you write. You write your module beside its testbench and run the
three GHDL commands \secref{c2:app:c} gives; the testbench prints its ``all checks passed'' note
when your module is right. At the repository's root:

\begin{shell}
make help                # List every target.
make build               # Build and simulate every example; check every exercise testbench.
make build MODULE=timer  # Build only the examples whose path contains "timer".
\end{shell}

\noindent Quartus Prime Lite and a DE0-CV board are needed for nothing in the book. If you have
them, Appendix~\ref{app:quartus} takes a design from VHDL onto the board.
